\lampiransection{Pengodean Tematik Lintas Sumber}
\label{app:pengodean-tematik-lintas-sumber}

Lampiran ini menyajikan ringkasan pengodean tematik lintas sumber yang digunakan untuk merangkum hubungan antara temuan lapangan, konsep pemasaran, faktor internal-eksternal, dan indikator peningkatan penjualan. Ringkasan ini membantu memperlihatkan bagaimana tema dari wawancara, angket pelanggan, observasi, dokumentasi, dan data pra-survei dibaca sebelum disintesiskan ke dalam analisis strategi.

\begingroup
\scriptsize
\begin{xltabular}{\textwidth}{|L{2.7cm}|Y|L{2.75cm}|L{3.25cm}|}
\caption{Ringkasan Pengodean Tematik Lintas Sumber}
\label{tab:pengodean-tematik-lintas-sumber-lampiran}\\
\toprule
\textbf{Tema Analisis} & \textbf{Sumber Kode Utama} & \textbf{Keterkaitan Konsep} & \textbf{Arah Pembacaan} \\ \hline
\midrule
\endfirsthead
\continuedtabletitle{4}
\toprule
\textbf{Tema Analisis} & \textbf{Sumber Kode Utama} & \textbf{Keterkaitan Konsep} & \textbf{Arah Pembacaan} \\ \hline
\midrule
\endhead
\continuedtablefooter{4}
\endfoot
\bottomrule
\endlastfoot
Cita rasa lauk kuat & K01.5, K05.2, K05.7, K04.3, K02.8, K02.9, K08.1, K08.4, K09.2--K09.4 & \textit{Product}, \textit{positioning}, nilai pelanggan & Dibaca sebagai kekuatan karena mendorong minat beli awal, pembelian ulang, dan peluang paket menu. \\ \hline
Persepsi harga GrabFood & K01.19, K04.5, K04.6, K02.10, K02.17, K08.1, K09.5 & \textit{Price}, \textit{promotion}, keputusan \textit{checkout} & Dibaca sebagai kelemahan atau ancaman ketika harga akhir dan margin tidak terkendali. \\ \hline
Akurasi packing & K01.14, K01.15, K01.16, K01.3, K01.4, K04.9, K05.8, K02.5, K02.15, K02.20 & \textit{Process}, \textit{people}, kualitas layanan & Dibaca sebagai kelemahan utama karena berkaitan langsung dengan item kurang, waktu tunggu, komplain, rating, dan pembelian ulang. \\ \hline
Update stok & K05.6, K01.8, K04.7, K08.1, K03.6 & \textit{Place}, \textit{process}, ketersediaan penawaran & Dibaca sebagai kelemahan yang dapat menghambat konversi karena pelanggan dapat membatalkan pesanan atau memilih kompetitor. \\ \hline
Tampilan toko digital & K05.1, K05.2, K04.7, K04.8, K02.11, K01.22 & \textit{Physical evidence}, pemasaran digital, kepercayaan pelanggan & Dibaca sebagai kekuatan yang perlu diperkuat sekaligus kelemahan pada bagian informasi paket dan standar tampilan layanan. \\ \hline
Segmen pelanggan dan kebutuhan makan praktis & K04.1, K04.2, K02.7, K02.17 & Segmentasi, \textit{targeting}, perilaku pelanggan OFD & Dibaca sebagai peluang untuk paket makan siang pekerja, paket mahasiswa, pelanggan sekitar, dan pelanggan yang sensitif terhadap ongkir. \\ \hline
Digitalisasi internal, sistem antrean, menu digital, dan \textit{stamp digital} & K02.1--K02.5, K02.11, K02.12, K02.16, K02.18--K02.21, K09.6, K09.10--K09.12 & \textit{Process}, loyalitas, pengendalian layanan & Dibaca sebagai kekuatan awal dan peluang perbaikan karena dapat mengurangi antrean, mempercepat pemesanan, dan membangun data pelanggan. \\ \hline
Reputasi kompetitor & K06.1--K06.5, K03.3--K03.8, K03.12, K03.15, K03.18 & Lingkungan mikro, persaingan, diferensiasi & Dibaca sebagai ancaman karena pelanggan mudah membandingkan merchant, tetapi juga sebagai peluang ketika kompetitor memiliki kelemahan proses yang serupa. \\ \hline
Indikator kinerja, AOV, kunjungan, dan pembelian ulang & K08.3, dokumentasi kinerja, K02.18--K02.21, K09.5, K09.10 & Indikator peningkatan penjualan dan evaluasi strategi & Dibaca sebagai konteks urgensi bahwa strategi perlu memperbaiki konversi, pengalaman layanan, dan loyalitas, bukan hanya menambah trafik. \\ \hline
\end{xltabular}
\tablesource[\textwidth]{Diolah peneliti berdasarkan koding wawancara, angket, observasi, dokumentasi, dan konteks pra-survei (2026).}
\endgroup
